\documentclass[../../../../SoftwareRequirementsSpecification.tex]{subfiles}
\begin{document}
% Richiede le macro definite in src/macros/useCaseMacros.tex
% (incluse nel preambolo di SoftwareRequirementsSpecification.tex)

\ucstart
\begin{xltabular}{\textwidth}{|l|X|}
  \hline
  \ucid{PT-13} \\ \hline
  \ucname{Associazione a una Palestra} \\ \hline
  \ucactor{PT} \\ \hline
  \ucdescription{Il PT si associa a una palestra in cui lavora,
    indicando la data di inizio del rapporto.} \\ \hline
  \ucprecond{Il PT deve essere autenticato. La palestra a cui il PT si
    vuole associare deve essere già registrata (vedi caso d'uso
    PT-12).} \\ \hline
  \ucpostcond{Il sistema crea l'associazione tra il PT e la palestra,
    con la data di inizio indicata (vedi Nota 1). L'associazione
    compare nell'elenco del caso d'uso PT-14.} \\ \hline
  \ucmainflow{
    \item Il PT apre la pagina di ricerca palestra.
    \item Il sistema mostra un campo di ricerca.
    \item Il PT inserisce il nome della palestra a cui vuole
      associarsi.
    \item Il sistema mostra l'elenco delle palestre che corrispondono
      alla ricerca, con nome e indirizzo.
    \item Il PT preme su una palestra dell'elenco.
    \item Il sistema mostra il form di associazione, con la data di
      inizio precompilata alla data odierna.
    \item Il PT conferma o modifica la data di inizio e preme il
      pulsante "Associati".
    \item Il sistema verifica i dati.
    \item Il sistema crea l'associazione tra il PT e la palestra.
    \item Il sistema mostra una pagina di conferma.
  } \\ \hline
  \ucaltflow{
    \textbf{4a. Nessuna palestra trovata} \newline
    - Il sistema mostra un messaggio ("Nessuna palestra trovata") e il
    pulsante "Registra nuova Palestra". \newline
    - Se il PT preme il pulsante, il flusso continua dal passo 3 del
    caso d'uso PT-12; al termine, il flusso riprende dal passo 6 con la
    palestra appena creata. \newline
    \textbf{8a. Associazione già attiva} \newline
    - Il sistema mostra un messaggio di errore ("Sei già associato a
    questa palestra") e non crea l'associazione. \newline
    - Il flusso termina. \newline
    \textbf{8b. Data di inizio non valida} \newline
    - Il sistema mostra un messaggio di errore ("La data di inizio non
    può essere nel passato"). \newline
    - Il flusso riprende dal passo 6.
  } \\ \hline
  \ucnote{
    \textbf{Nota 1:} l'associazione resta attiva finché il PT non la
    termina con il caso d'uso PT-15, che ne imposta la data di fine.
    \newline
    Il PT può essere associato a più palestre contemporaneamente: ogni
    associazione è indipendente dalle altre.
  } \\ \hline
\end{xltabular}
\end{document}
